\section{Membaca dan Memeriksa Keluaran Antara}
\label{sec:ch03-runtime}

Tabel penelusuran hanya berguna jika isinya dapat diperiksa. Bagian
ini membahas tiga keluaran antara yang wajib muncul: aliran token,
pohon sintaks, dan kode antara.

\subsection{Aliran Token}

Aliran token adalah daftar terurut yang setiap elemennya memuat jenis
token, lexeme, dan lokasi. Untuk program pada bagian sebelumnya,
awal alirannya adalah:

\begin{center}
\small
\code{INT} \code{ID(main)} \code{LPAREN} \code{RPAREN}
\code{LBRACE} \code{INT} \code{ID(a)} \code{ASSIGN}
\code{INT\_LIT(5)} \code{SEMI} \code{INT} \code{ID(b)}
\code{ASSIGN} \code{INT\_LIT(7)} \code{SEMI} \code{PRINT}
\code{LPAREN} \code{ID(a)} \code{PLUS} \code{ID(b)}
\code{RPAREN} \code{SEMI} \code{RETURN} \code{INT\_LIT(0)}
\code{SEMI} \code{RBRACE} \code{EOF}.
\end{center}

Saat memeriksa aliran ini, cocokkan setiap token dengan lexeme pada
baris sumber. Kesalahan urutan atau token yang hilang menunjukkan
masalah pada aturan leksikal.

\subsection{Pohon Sintaks}

Pohon sintaks menunjukkan bagaimana token disusun menjadi struktur.
Representasi ringkas untuk program yang sama adalah:

\begin{verbatim}
Program
  DeklarasiFungsi(main)
    Blok
      Deklarasi(a, int, 5)
      Deklarasi(b, int, 7)
      Pernyataan(print)
        Ekspresi(+)
          Variabel(a)
          Variabel(b)
      Pernyataan(return)
        Bilangan(0)
\end{verbatim}

Setiap simpul anak harus berasal dari token yang benar. Jika sebuah
ekspresi muncul sebagai anak dari deklarasi, padahal seharusnya menjadi
anak pernyataan, pohon tersebut tidak sesuai dengan tata bahasa.

\subsection{Kode Antara}

Kode antara yang dipakai pada buku ini adalah instruksi tiga alamat.
Untuk program contoh, hasilnya dapat dituliskan sebagai:

\begin{center}
\small
\code{t1 = a + b}; \code{param t1}; \code{call print, 1};
\code{t2 = 0}; \code{return t2}.
\end{center}

Instruksi tiga alamat memisahkan operasi kompleks menjadi langkah
sederhana dengan nama sementara. Bentuk ini memudahkan pemeriksaan
urutan operasi dan menjadi titik awal pembangkitan kode target.

\subsection{Ciri Keluaran yang Benar}

\begin{table}[H]
\centering
\caption{Ciri keluaran antara yang dapat diperiksa}
\label{tab:ch03-output-check}
\begin{tabularx}{\textwidth}{@{}lXX@{}}
\toprule
\textbf{Keluaran} & \textbf{Ciri yang diperiksa} & \textbf{Contoh dari program} \\
\midrule
Aliran token & Jenis, lexeme, dan lokasi setiap token &
\code{ID(a)} pada baris 2 \\
Pohon sintaks & Hubungan induk--anak sesuai tata bahasa &
\code{Ekspresi(+)} menjadi anak \code{print} \\
AST teranotasi & Tipe dan referensi simbol pada setiap simpul &
\code{a + b} bertipe \code{int} \\
Kode antara & Operasi sederhana dan urutan yang konsisten &
\code{t1 = a + b}; \code{param t1} \\
\bottomrule
\end{tabularx}
\end{table}

\subsection{Langkah Mengisi Tabel}

\begin{enumerate}[label=\arabic*.,leftmargin=*]
  \item Pilih satu program C sederhana, maksimal sepuluh baris.
  \item Tandai token penting beserta baris dan kolomnya.
  \item Susun pohon sintaks dari aliran token.
  \item Tambahkan tipe dan referensi simbol pada simpul yang relevan.
  \item Tuliskan kode antara tiga alamat.
  \item Masukkan hasil setiap langkah ke tabel penelusuran dan periksa
        konsistensinya dengan baris program sumber.
\end{enumerate}

\begin{warningbox}
Jika pohon sintaks tidak konsisten dengan aliran token, jangan
mengarang keluaran. Periksa kembali aturan tata bahasa atau token yang
salah, lalu perbaiki tabel dari awal.
\end{warningbox}
